% Terminal note of the homily: everything the speech itself must not carry ---
% the audience, the spoken word count and the pace it implies, the relation to
% the three interpretations of the expansive study, the connections that are
% the homily's own, the translation each quoted text is given in, and the exact
% loci --- followed by the References. Nothing here is read aloud.

\section*{Note on Sources and Delivery}
\runninghead{Sources and Delivery}

{\footnotesize

\textbf{Audience and occasion.} An adult parish assembly, at the Mass
\latin{Dominica vigesima quarta et ultima post Pentecosten}, II classis, of
the 1962 \work{Missale Romanum}, editio typica, pp.~422--424,
nos.~1769--1778, said on Sunday 22~November 2026, the last Sunday after
Pentecost. The speech presumes that it follows the Gospel, and it speaks of
the Secret and the Communion as still to come. It uses the Introit, the
Epistle, the Gradual, the Gospel, the Secret and the Communion; it does not
use the Collect, the Alleluia, the Offertory or the Postcommunion. No prayer
is recited.

\textbf{Length and pace.} The spoken body runs to about 1,340 words (1,342 by
a plain count with comments and markup removed), which is about 10.3 to 11.2
minutes at an unhurried 120 to 130 words a minute, by arithmetic alone; no
one has delivered or timed it. It was read through in full, silently, for
sense and ease of speech.

\textbf{Relation to the reviewed interpretations.} The homily follows the
third of the expansive study's three interpretations, \textit{My Words Shall
Not Pass: What God Has Said Is as Good as Done}: Chrysostom on ``this
generation'' and on the words that do not pass, Gregory on heaven and earth
and on gladness at the world's end, Augustine on the Gradual's past tense,
Chrysostom on the Epistle's thanks, Jerome on the Introit's call to prayer,
and Bede on the condition of the Communion's promise. It takes the literal
sense of the commands of flight from the second interpretation, which sets
it first: Chrysostom on the Roman war and Bede on the flight to Pella. The
first interpretation, on the body and the eagles and the captivity brought
back, is left to the studies, and so are the spiritual readings of the
flight (Hilary's housetop, Jerome's perverse doctrine), the other readings
of the fig tree, Jerome's and Hilary's answers on ``this generation,'' and
Chrysostom's identification of the abomination (\work{Hom.\ in Mt.}~75.2).

\textbf{Connections that are the homily's own.} These are drawn in the
preacher's voice and credited to no Father. The gloss of ``abomination'' as
a thing God abhors, which leaves the holy place desolate, explains the
English word, and naming Vespasian as the Roman general identifies him for
the hearer. That the warning was obeyed rests on Bede's account (\latin{Haec
juxta litteram facta esse constat}); that the fulfilled first part of the
Gospel is a ground for believing the rest is the preacher's application of
Chrysostom's sentence on v.~35, which concerns the Lord's sayings in general
(77.1), and agrees with the second interpretation's statement that the history
shows the command obeyed. Chrysostom explains the lightning of v.~27 as a
coming that needs no herald (76.3); that the Lord gives no date rests on
v.~36, before which the Gospel stops. That everyone who believes belongs to
Chrysostom's ``generation of the believers'' and that the promise is made to
the hearers is the preacher's. Augustine expounds Ps 43:8 in his own lemma,
\latin{Salvos enim fecisti nos}, where the Gradual has \latin{Liberásti nos}; the application to the Gradual is
the preacher's, as it is the study's. Chrysostom's ``causes'' are applied to
the hearers' forgiveness and place in the kingdom by the preacher. Jerome
comments on Jer 29:12 and not on the Mass; the step from his sentence to
continued prayer for what God has promised, and the remark that the Communion
puts promise and prayer in one sentence, are the preacher's. Bede comments on
Mk 11:25 and does not mention the Mass; the response, to forgive before
Communion and to ask forgiveness believing that it will be received, applies
Mk 11:24--25 to the hearers. Gregory's sentence on gladness at the world's end
comments on Lk 21:28 and does not mention the Secret; the closing joins them
in the preacher's voice.

\textbf{Sources, translations and loci.} Each source is listed once, with the
translation the speech uses; ``paraphrased'' means told in the speech's own
words. \emph{The Mass}, nos.~1769--1778. The Introit antiphon (Jer
29:11--12, 14; p.~736), the Gradual (Ps 43:8; pp.~737--738), the Secret
(p.~740) and the Communion (Mk 11:24; p.~740) are quoted in the English of
the Lasance \work{New Roman Missal} (1945), a historical study aid and not
an approved liturgical translation of the 1962 Missal; the Secret is quoted
from ``that, being delivered,'' and its petition to turn all our hearts to God
is paraphrased in its words.
\emph{Scripture} is quoted in the Douay--Rheims, Challoner revision, at Mt
24:15, 23, 27, 34, 35 and 36, Ps 23:6, Col 1:11 and 13 and Mk 11:25; Mt
24:1--3, 16--20, 24, 30--33, Ps 43:10, Jer 29:10 and Col 1:10, 14 are told in
the speech's own words. \emph{St John Chrysostom}, in the Nicene and Post-Nicene
English: \work{Hom.\ in Mt.}~76.1, the sabbath and Vespasian, paraphrased;
76.3, the lightning that ``needs not a herald'' and shows itself everywhere
at once, paraphrased; 77.1, ``from what is past, let him believe also the
things to come,'' ``not of the generation then living, but of that of the
believers'' and ``over this generation shall nothing prevail, not famine, not
pestilence,'' quoted; \work{Hom.\ in Col.}~2, on Col 1:12, that Paul
``mentions the things which were already theirs, for these are the causes of
the other,'' paraphrased closely. \emph{St Bede}, in the expansive study's
English of the Migne Latin: \work{In Marci evangelium expositio}~IV, on
13:14, the Christians of the province \latin{oraculo admoniti} withdrawing
across the Jordan to Pella, paraphrased; III, on 11:25, \latin{Qui vero tantae
fastigium perfectionis necdum conscendere queunt, postulent sibi peccata
dimitti \dots\ et absque dubio impetrabunt quod petunt, si tamen in se
peccantibus aliis primo ipsi dimittunt}, paraphrased closely. \emph{St
Gregory the Great}, in the expansive study's English of the Migne Latin:
\work{Homiliae in Evangelia}~I.1.4, on Lk 21:33, \latin{Nihil enim in rerum
corporalium natura coelo et terra durabilius, et nihil in rerum natura tam
velociter quam sermo transit}, with the sentence on words that are not
complete until they pass, and \latin{sine mutabilitate manentes sententias
exprimit meus sermo qui transit}, paraphrased closely; I.1.3, \latin{Qui ergo
Deum diligunt ex mundi fine gaudere atque hilarescere jubentur, quia
videlicet eum quem amant mox inveniunt, dum transit is quem non amaverunt},
paraphrased closely. \emph{St Augustine}, in the expansive study's English of
the Migne Latin: \work{Enarrationes in Psalmos}~43.8, \latin{Et hoc figura
praeteriti de futuro dicitur: sed ideo tanquam praeteritum dicitur, quia tam
certum est quasi factum sit}, \latin{Nobis enim ea quae praeterierunt, certa
sunt; quae futura, incerta sunt} and \latin{Da prophetam cui tam certum sit
futurum quam tibi praeteritum}, paraphrased closely, with his sentence on what
cannot not have happened. \emph{St Jerome}, in the expansive study's English
of the Migne Latin: \work{In Hieremiam}~V, on 29:10--14, PL~24 col.~893,
\latin{Certe absque invocatione et oratione captivorum, poterat implere
Dominus quod promiserat; sed hortatur eos ad preces, ut quod promissum est,
mereantur accipere}, paraphrased closely.

\textbf{Text checked.} The Lasance English was taken from
\texttt{propers/verified.md}, and the Douay--Rheims from the registered
Challoner verse texts of Matthew, Mark, Psalms, Jeremias and Colossians.
Chrysostom was re-read in the registered Nicene and Post-Nicene texts of
vols.~10 and 13, Gregory and Bede in their registered Wikisource Latin texts,
and Augustine in the registered Migne Latin passage of 43.8. Jerome's
sentence on Jer 29 was not re-read for the homily: it is given as the
research record read it on the page image of PL~24, col.~893.

\par}

\Needspace{14\baselineskip}
\section*{References}
\runninghead{References}

{\footnotesize
\begin{itemize}\setlength{\itemsep}{0.15em}\setlength{\parskip}{0pt}
\item \work{Missale Romanum ex decreto Sacrosancti Concilii Tridentini restitutum, Summorum Pontificum cura recognitum}, editio typica (Typis Polyglottis Vaticanis, 1962), CMAA facsimile, Proprium de Tempore, pp.~422--424, nos.~1769--1778.
\item Lasance, F.~X., and Walsh, F.~A., \work{The New Roman Missal in Latin and English}, revised edition (New York: Benziger Brothers, 1945), pp.~736--740 (Introit, Gradual, Secret, Communion).
\item \work{The Holy Bible Translated from the Latin Vulgate} (Douay--Rheims, Challoner revision), Project Gutenberg text: Ps 23:6; 43:10; Jer 29:10; Mt 24:1--3, 15--36; Mk 11:25; Col 1:10--14.
\item St Augustine, \work{Enarrationes in Psalmos} 43.8, Migne Latin text (Wikisource).
\item St Bede, \work{In Marci evangelium expositio} III, on 11:25, and IV, on 13:14, Migne Latin text (Wikisource).
\item St John Chrysostom, \work{Homilies on Matthew} 76.1, 76.3 and 77.1, Nicene and Post-Nicene Fathers, 1st ser., vol.~10; \work{Homilies on Colossians} 2, Nicene and Post-Nicene Fathers, 1st ser., vol.~13.
\item St Gregory the Great, \work{Homiliae in Evangelia} I.1, nn.~3--4, Migne Latin text (Wikisource).
\item St Jerome, \work{In Hieremiam} V, on 29:10--14, PL~24 (Paris, 1865), col.~893.
\end{itemize}
\par}
